\documentclass[12pt]{report}
\usepackage[spanish]{babel}
\usepackage[utf8]{inputenc}
\usepackage[T1]{fontenc}
\usepackage{lmodern}
\usepackage{geometry}
\geometry{a4paper,margin=2.5cm}
\usepackage{hyperref}
\hypersetup{colorlinks=true,linkcolor=blue,urlcolor=blue}


\title{Avances}
\author{Javier Zavala Ponce}
\date{\today}

\begin{document}
	
	% --- Portada ---
	\begin{titlepage}
		\centering
		\vspace*{2cm}
		{\LARGE ESIME Culhuacán\par}
		\vspace{1cm}
		{\Large IPN\par}
		\vspace{2cm}
		{\bfseries\Huge Reporte de actividades 001\par}
		\vspace{2cm}
		{\large Javier Zavala Ponce\par}
		{\large Control Automático\par}
		\vfill
		{\large \today\par}
	\end{titlepage}
	
	% --- Índice (tabla de contenidos) ---
	\tableofcontents
	\newpage
	
	% === Tema 1 ===
	\chapter{Sobre Libro Thinking in Systems}
	\section{Resumen}
Lo percibo como una especie de guía para entender por qué los problemas complejos (como la pobreza, el desempleo o el daño ambiental) persisten a pesar de las medidas que se han ido tomando. Donella Meadows, la autora, invita a dejar de buscar culpables externos y a observar la estructura interna de los sistemas al rededor de nosotros.


\begin{itemize}
	\item Un sistema es un conjunto de elementos, interconexiones y un propósito. 
	
	\item Cambiar los elementos afecta; pero cambiar las interconexiones o el propósito transforma todo.
	

\item Los sistemas tienen existencias (stocks) que cambian lentamente gracias a flujos (entradas y salidas). Esto genera inercia: las cosas no cambian de la noche a la mañana.

\item El comportamiento del sistema contiene lazos de retroalimentación:

\item Los sistemas nos sorprenden porque actuamos sin ver las demoras, los bucles ocultos o las metas contradictorias.

\item Palancas de cambio: lo más poderoso no es meter más dinero o gente, sino cambiar las reglas, los flujos de información o el propósito del sistema.

\item Conclusión: Para mejorar un sistema, hay que entender su estructura, respetar sus demoras y evitar soluciones mágicas.

\end{itemize}
	
	\section{Comentarios}
Lo que mas llamo mi atencion es la raigambre ideologica de Meadows (...)
encontre referencias como Wiener (quien es considerado el fundador de la cibernetica), 
y tambien Bertalanffy.

Creo que el motivo detras de las imprecisiones del control de algunos sistemas, podria ubicarse en el modelo, e inclusive en el paradigma ideologico y la base epistemologica con
que el observador realiza un modelo. 
	
	% === Tema 2 ===
	\chapter{Tema 2: Ecuaciones Diferenciales}
	\section{Avance}
	Tengo copia del Zill. Reconozco algunas lagunas de conocimiento y estoy en proceso de desenpolvar
	libros y conceptos.
	
	\section{Comentario	}
He comenzado a utilizar julia como lenguaje para reforzar conceptos en computadora
	
	% === Tema 3 ===
	\chapter{Tema 3: Libro de control: Bolton}
	\section{Avances}

Cero avance con este libro
sin embargo, consegui tambien el Ogata y el Kuo para referencias posteriores

	\section{Comentarios}
El enfasis de la sesion pasada estuvo en la notacion, lo cual tengo presente para
compatibilizar mis reportes.	
\end{document}
